% ============================================================
%  Paper 3 — Aroha Labs
%  "Tiered Publisher Verification for Decentralised
%   AI Agent Registries: A Security Analysis"
%
%  Target venues: IEEE S&P 2027 / ACM CCS 2027 / WWW 2027
%  arXiv categories: cs.CR, cs.NI, cs.AI
% ============================================================

\documentclass[conference]{IEEEtran}
\IEEEoverridecommandlockouts

\usepackage{amsmath,amssymb,amsthm}
\usepackage{mathtools}
\usepackage{algorithm}
\usepackage{algpseudocode}
\usepackage{graphicx}
\usepackage{booktabs}
\usepackage{hyperref}
\usepackage{xcolor}
\usepackage{cite}
\usepackage{bm}
\usepackage{tabularx}

\newtheorem{theorem}{Theorem}
\newtheorem{lemma}[theorem]{Lemma}
\newtheorem{proposition}[theorem]{Proposition}
\newtheorem{corollary}[theorem]{Corollary}
\newtheorem{definition}{Definition}
\newtheorem{remark}{Remark}
\newtheorem{attack}{Attack}

\DeclareMathOperator{\negl}{\mathsf{negl}}
\DeclareMathOperator{\E}{\mathbb{E}}
\newcommand{\Adv}{\mathcal{A}}
\newcommand{\pk}{\mathsf{pk}}
\newcommand{\sk}{\mathsf{sk}}
\newcommand{\sig}{\sigma}
\newcommand{\Sign}{\mathsf{Sign}}
\newcommand{\Vrfy}{\mathsf{Vrfy}}
\newcommand{\Gen}{\mathsf{Gen}}
\newcommand{\Fetch}{\mathsf{Fetch}}

\begin{document}

\title{Tiered Publisher Verification for Decentralised\\
       AI Agent Registries: A Security Analysis}

\author{
  \IEEEauthorblockN{Aroha Labs Research}
  \IEEEauthorblockA{%
    \textit{Aroha Protocol}\\
    aroha-labs.com
  }
}

\maketitle

% ─────────────────────────────────────────────────────────────
\begin{abstract}
Decentralised AI agent registries face a fundamental
impersonation problem: any participant can claim to be any
organisation when registering an agent, since the registry
has no mechanism to verify publisher identity at registration
time.  We present a three-tier publisher verification
architecture—community, domain, and org-signed—that
progressively raises the cost and cryptographic strength
of publisher identity claims. The domain tier leverages DNS
control as a trust signal via a well-known URI binding
\cite{rfc5785}, while the org-signed tier uses Ed25519
keypairs to produce cryptographically verifiable claims
about the publishing organisation.  We additionally introduce
a \emph{discovery cooldown} mechanism—a 48-hour period during
which newly registered agents are excluded from public
discovery lists but remain directly resolvable—as an
anti-impersonation and anti-spam gate.
We formalise the threat model, prove security properties for
each tier, characterise the attack cost for each threat, and
provide a comprehensive comparison against alternative
approaches (X.509 WebPKI, Sigstore, DNSSEC).  The system is
deployed in production on the Aroha agent registry, handling
real-time publisher verification queries.
\end{abstract}

\begin{IEEEkeywords}
publisher verification, AI agent registry, decentralised
identity, Ed25519, well-known URI, Sybil resistance,
impersonation, discovery cooldown, DID
\end{IEEEkeywords}

% ─────────────────────────────────────────────────────────────
\section{Introduction}
\label{sec:intro}

Public agent registries—indexes in which organisations
publish AI agents that other systems can discover and
invoke—are emerging as critical infrastructure for
autonomous multi-agent AI systems.  Unlike software package
registries (npm, PyPI), agent registries describe autonomous
entities that can spend money, access external systems, and
make consequential decisions. The identity of the publisher
is therefore a critical trust signal for potential users.

The challenge is that a registry open to all participants
cannot rely on any centralised identity authority.  Any
actor can claim to be OpenAI, Anthropic, or any other
trusted organisation simply by setting the \texttt{orgName}
field in a registration request.  This creates three distinct
threat classes:

\begin{enumerate}
  \item \textbf{Impersonation}: A malicious actor registers
    an agent claiming to be from a trusted organisation to
    lend credibility to a harmful agent.
  \item \textbf{Sybil flooding}: An adversary registers
    large numbers of low-effort agents to dominate discovery
    results and crowd out legitimate agents.
  \item \textbf{Reputation laundering}: An agent with a
    declining reputation is re-registered under a new
    identity to escape its history.
\end{enumerate}

We design a verification architecture that raises the
attack cost for each threat class proportionally to the
claimed trust level, without requiring a centralised
certificate authority or blockchain.

Our contributions:
\begin{enumerate}
  \item A formal threat model for open AI agent registries
    with three attacker profiles (\S\ref{sec:threat}).
  \item A three-tier publisher verification architecture
    with formal security proofs for each tier
    (\S\ref{sec:tiers}).
  \item A discovery cooldown mechanism with formal analysis
    of its anti-Sybil and anti-impersonation properties
    (\S\ref{sec:cooldown}).
  \item An adversarial cost analysis comparing all three
    tiers (\S\ref{sec:cost_analysis}).
  \item A comparison to X.509/WebPKI, DNSSEC, and Sigstore
    (\S\ref{sec:comparison}).
\end{enumerate}

% ─────────────────────────────────────────────────────────────
\section{Background}
\label{sec:background}

\subsection{Decentralised Identifiers}

The Aroha protocol uses Decentralised Identifiers (DIDs)
\cite{w3c2022did} as agent identity anchors.  Each agent
has a DID of the form:
\begin{equation}
  d = \texttt{did:aroha:studio:}h
  \label{eq:did}
\end{equation}
where $h = \mathsf{Keccak256}(\text{utf8}(d))$ is the
didHash—a 256-bit Keccak hash \cite{keccak2015} of the DID
string itself (used as a stable content-addressed key).
The didHash is computed deterministically and serves as the
primary key in the registry database.

\subsection{Content-Addressed Manifests}

Each agent's capabilities are described in a manifest file
pinned to IPFS \cite{ipfs2014whitepaper} and referenced by
its CID (Content Identifier).  The CID is the SHA-256 hash
of the manifest content encoded in the IPLD raw block
format.  This makes the manifest content-addressed and
tamper-evident: any modification to the manifest invalidates
the CID stored in the registry, detectable by the manifest
integrity signal (\S~Paper~1).

\subsection{Well-Known URIs}

RFC 5785 \cite{rfc5785} defines a convention for hosting
per-domain metadata at \texttt{/.well-known/<name>}. The
ACME protocol \cite{acme2019rfc} uses this mechanism to
prove domain control for TLS certificate issuance—the same
security property we leverage for domain-tier verification.

% ─────────────────────────────────────────────────────────────
\section{Threat Model}
\label{sec:threat}

\subsection{System Components}

We model the registry as consisting of:
\begin{itemize}
  \item \textbf{Registry server}: stores and serves agent
    records; implements the verification API.
  \item \textbf{Agents}: participants who register agents.
  \item \textbf{Consumers}: applications that query the
    registry to discover agents.
  \item \textbf{DNS}: the global domain name system
    (assumed honest but not authenticated without DNSSEC).
\end{itemize}

\subsection{Attacker Profiles}

\begin{definition}[Attacker Profile]
We define three attacker profiles with increasing capability:
\begin{enumerate}
  \item $\Adv_1$ (Basic): API key holder; can register
    agents with arbitrary metadata; cannot forge Ed25519
    signatures; cannot control DNS records for domains they
    do not own.
  \item $\Adv_2$ (Domain): All capabilities of $\Adv_1$
    plus: can register arbitrary domain names and set their
    DNS records; cannot compromise the private keys of
    legitimate organisations.
  \item $\Adv_3$ (Advanced): All capabilities of $\Adv_2$
    plus: can perform DNS cache poisoning attacks with
    probability $p_{\text{dns}} \ll 1$; cannot break
    Ed25519 (assumed EUF-CMA secure).
\end{enumerate}
\end{definition}

\subsection{Security Goals}

\begin{definition}[Security Goals]
The registry aims to satisfy:
\begin{enumerate}
  \item \textbf{Impersonation resistance}: $\Adv_1$ cannot
    register an agent as published by organisation $O$
    without $O$'s cooperation.
  \item \textbf{Domain integrity}: An agent claiming
    domain $d$ at tier $\geq 1$ is only registered if the
    agent's DID appears in the well-known resource at $d$.
  \item \textbf{Org-key integrity}: An agent at tier 2
    is only registered if the org's private key was used
    to sign the registration payload.
  \item \textbf{Sybil resistance}: An adversary cannot
    flood the discovery index with many low-effort agents
    without paying proportional cost per agent.
\end{enumerate}
\end{definition}

% ─────────────────────────────────────────────────────────────
\section{Three-Tier Verification Architecture}
\label{sec:tiers}

\begin{definition}[Verification Tier]
Each registered agent is assigned one of three tiers
$T \in \{T_0, T_1, T_2\}$, assigned at registration time
or upgraded via explicit verification calls:
\begin{itemize}
  \item $T_0$ (Community): Default tier. No cryptographic
    claims about publisher identity.
  \item $T_1$ (Domain): The agent's DID has been verified
    as appearing in the DNS-controlled well-known resource
    for a claimed domain $d$.
  \item $T_2$ (Org-signed): The registration payload was
    signed by an Ed25519 key associated with the claimed
    organisation via a well-known org-key binding.
\end{itemize}
\end{definition}

\subsection{Tier 0: Community}

All registered agents start at $T_0$.  The community tier
provides no guarantees about publisher identity.  It is
suitable for:
\begin{itemize}
  \item Individual developers who have not set up domain
    infrastructure.
  \item Early-stage agents not yet ready for domain binding.
  \item Cases where the consuming application does not
    require strong publisher guarantees.
\end{itemize}

The cost of registration at $T_0$ is the cost of an API key:
$\mathcal{C}_0 = c_{\text{key}}$. This is non-zero (API keys
require a verified email account), providing minimal Sybil
resistance.

\subsection{Tier 1: Domain Verification}

Domain-tier verification proceeds as follows:

\begin{definition}[Domain Verification Protocol]
\label{def:domain_verify}
To achieve $T_1$ for domain $d$ and agent DID $\delta$:
\begin{enumerate}
  \item The publisher hosts a JSON file at:
    \[
      \texttt{https://}d\texttt{/.well-known/aroha-did.json}
    \]
    with content $\{\texttt{"did"}: \delta\}$.
  \item The publisher calls:
    \[
      \texttt{POST /v1/agents/}\hat{\delta}
        \texttt{/verify-domain}
    \]
    with body $\{\texttt{"domain"}: d\}$, authenticated
    by their API key.
  \item The registry fetches the well-known resource and
    verifies:
    {\small\begin{align}
      &\Fetch\!\left(d,\;\texttt{/.well-known/aroha-did.json}
        \right)\texttt{.did} = \delta
        \label{eq:domain_check}
    \end{align}}
  \item On success, the registry sets the agent's tier to
    $T_1$ and records $d$ as the verified domain.
\end{enumerate}
\end{definition}

\begin{proposition}[Domain Tier Security]
\label{prop:domain}
Against $\Adv_1$, domain verification achieves
impersonation resistance for any domain $d$ that $\Adv_1$
does not control.
\end{proposition}
\begin{proof}
For $\Adv_1$ to register an agent as verified for domain
$d$, they must produce a well-known file at $d$ containing
$\delta_{\Adv}$. But $\Adv_1$ cannot modify DNS records for
$d$ (by assumption). Therefore the registry's fetch of
$d$\texttt{/.well-known/aroha-did.json} returns the
legitimate domain holder's content, which does not contain
$\delta_{\Adv}$, causing the check \eqref{eq:domain_check}
to fail.
\end{proof}

\begin{remark}[DNS Limitations]
$\Adv_2$ can achieve $T_1$ for any domain they register.
This means $T_1$ does not prevent an adversary from
claiming \texttt{openai-agents.com} instead of
\texttt{openai.com}. The domain tier proves control of the
specific domain, not legitimacy of the domain name.
Homograph attacks (e.g., using Punycode to register
look-alike domains) remain possible; consumers should
verify the full domain string, not just its visual
appearance.
\end{remark}

\subsection{Tier 2: Org-Signed Verification}

The org-signed tier introduces an Ed25519 keypair
associated with an organisation, binding the key to a
domain via a second well-known resource.

\begin{definition}[Org-Key Registration Protocol]
\label{def:org_key}
An organisation $O$ with domain $d$ registers an org key
$\pk_{\text{org}}$ by:
\begin{enumerate}
  \item Generating Ed25519 keypair $(\sk_{\text{org}},
    \pk_{\text{org}}) \leftarrow \Gen(1^\lambda)$.
  \item Hosting at $d$\texttt{/.well-known/aroha-org.json}:
    \[
      \{\texttt{"orgPublicKeyB64"}: \pk_{\text{org}}\}
    \]
  \item Calling \texttt{PUT /v1/auth/org-key} with
    $\{\texttt{"domain"}: d,\;
       \texttt{"orgPublicKeyB64"}: \pk_{\text{org}}\}$.
  \item The registry verifies the well-known binding:
    {\small\begin{align}
      &\Fetch\!\left(d,\;
         \texttt{/.well-known/aroha-org.json}\right)\nonumber\\
      &\quad\texttt{.orgPublicKeyB64} = \pk_{\text{org}}
        \label{eq:org_key_check}
    \end{align}}
    and stores $\pk_{\text{org}}$ associated with the API key.
\end{enumerate}
\end{definition}

\begin{definition}[Org-Signed Registration Protocol]
\label{def:org_signed}
To register an agent at $T_2$:
\begin{enumerate}
  \item The publisher computes:
    {\small\begin{align}
      p = \mathsf{Canon}\!\left(
        \bigl\{
          &\texttt{"action"}: \texttt{"register-agent"},\nonumber\\
          &\texttt{"did"}: \delta,\;
           \texttt{"manifestCID"}: m,\nonumber\\
          &\texttt{"publicKeyB64"}: k,\nonumber\\
          &\texttt{"registeredAt"}: t_r
        \bigr\}
      \right)
      \label{eq:org_payload}
    \end{align}}
    (keys in ASCII lexicographic order).
  \item The publisher computes $\sig_r = \Sign(\sk_{\text{org}},
    p)$.
  \item The registration request includes $\sig_r$ and $t_r$.
  \item The registry verifies:
    \begin{align}
      &\Vrfy(\pk_{\text{org}}, p, \sig_r) = 1
        \label{eq:org_verify} \\
      &|t_r - t_{\text{now}}| \leq \delta_{\text{fresh}} = 5\text{ min}
        \label{eq:org_fresh}
    \end{align}
    and on success, sets tier to $T_2$ and stores $\sig_r$.
\end{enumerate}
\end{definition}

\begin{theorem}[Org-Signed Security]
\label{thm:org_signed}
Against $\Adv_3$, the org-signed tier achieves
impersonation resistance for any organisation $O$ with
domain $d$ as long as $\sk_{\text{org}}$ is not compromised.
Formally, no PPT adversary can register an agent at $T_2$
for $O$ without knowledge of $\sk_{\text{org}}$ with
probability greater than $\max(\negl(\lambda),
p_{\text{dns}})$.
\end{theorem}
\begin{proof}
We split the proof into two cases.

\textbf{Case 1} ($\Adv_3$ does not mount a DNS attack).
For the registration to succeed, \eqref{eq:org_verify}
must hold: $\Vrfy(\pk_{\text{org}}, p, \sig_r) = 1$.
The payload $p$ encodes the adversary's DID $\delta_{\Adv}$
(different from any legitimate DID registered to $O$). By
EUF-CMA security of Ed25519, no PPT adversary can produce
a valid signature $\sig_r$ on $p$ without $\sk_{\text{org}}$
with probability greater than $\negl(\lambda)$.

\textbf{Case 2} ($\Adv_3$ mounts a DNS cache poisoning attack
to substitute $\pk_{\text{org}}$ with $\pk_{\Adv}$ in the
org-key check).  The adversary succeeds in the DNS attack
with probability $p_{\text{dns}}$. If the attack succeeds,
$\Adv_3$ can register $\pk_{\Adv}$ as the org key and then
sign registrations with the corresponding $\sk_{\Adv}$.
This attack succeeds with probability $p_{\text{dns}}$.
Under DNSSEC-protected domains, $p_{\text{dns}} \approx 0$;
under plain DNS, $p_{\text{dns}}$ is bounded by the cache
poisoning success rate (historically small but non-zero
\cite{acme2019rfc}).

Combining, the success probability is at most
$\max(\negl(\lambda), p_{\text{dns}})$.
\end{proof}

\begin{remark}[Limitation: Legal Identity]
\label{rem:legal}
$T_2$ proves possession of the private key corresponding to
a published org key, and proves that the key is linked to
a specific domain via DNS. It does not prove the legal
identity of the key holder. An adversary who registers
\texttt{anthropic-labs.io} and publishes an org key there
can register agents at $T_2$ for that domain. The gap
between domain control and legal identity requires out-of-band
vetting; Sigstore \cite{sigstore2021} with OIDC attestations
from cloud identity providers is the most promising
cryptographic approach to closing this gap in future work.
\end{remark}

% ─────────────────────────────────────────────────────────────
\section{Discovery Cooldown}
\label{sec:cooldown}

\subsection{Motivation}

Even with tiered verification, a patient attacker can
register a plausible agent, wait for it to accumulate
reputation via manufactured signals, and then use the
registry listing to contact victims. The critical window is
the first 48 hours after registration: this is when the
agent appears fresh in the registry, before any quality
signals have been submitted and validated.

We introduce the \emph{discovery cooldown}: a time window
$\tau_c$ after registration during which an agent is
excluded from the public discovery index but remains
directly resolvable by its didHash.

\subsection{Formal Definition}

\begin{definition}[Discovery Cooldown]
\label{def:cooldown}
At registration time $t_{\text{reg}}$, each agent is
assigned:
\begin{equation}
  t_{\text{elig}} = t_{\text{reg}} + \tau_c
  \label{eq:eligibility}
\end{equation}
The discovery index returns agent $a$ in response to list
queries if and only if:
\begin{equation}
  t_{\text{now}} \geq t_{\text{elig}}
  \;\vee\;
  t_{\text{elig}} = \bot
  \label{eq:discovery_gate}
\end{equation}
where $t_{\text{elig}} = \bot$ indicates no cooldown
(used for legacy records). Direct resolution by didHash
$\hat{\delta}$ is always available regardless of
$t_{\text{elig}}$.

The deployed value is $\tau_c = 48\text{ hours}$.
\end{definition}

\subsection{Anti-Impersonation Analysis}

\begin{definition}[Impersonation Campaign Model]
An impersonation campaign consists of:
\begin{enumerate}
  \item Registration of agent $\hat{\delta}$ at time $t_0$
    claiming identity of organisation $O$.
  \item Discovery and engagement with victim applications
    in time window $[t_{\text{elig}}, t_{\text{detect}}]$.
  \item $t_{\text{detect}}$ is the time at which $O$
    detects and requests removal of $\hat{\delta}$.
\end{enumerate}
\end{definition}

\begin{theorem}[Cooldown Effectiveness]
\label{thm:cooldown}
Let $T_{\text{detect}}$ be the expected detection time for
a legitimate organisation $O$ monitoring its name in the
registry.  Under the discovery cooldown with $\tau_c =
48\text{ h}$:
\begin{enumerate}
  \item If $T_{\text{detect}} \leq \tau_c$, the
    impersonating agent has zero exposure in the public
    discovery index before removal.
  \item If $T_{\text{detect}} > \tau_c$, the exposure
    window is $T_{\text{detect}} - \tau_c < T_{\text{detect}}$.
\end{enumerate}
\end{theorem}
\begin{proof}
(1) If $O$ detects the registration within $\tau_c$ hours,
it can request removal before $t_{\text{elig}}$ is reached.
By \eqref{eq:discovery_gate}, the agent never appears in
discovery results since $t_{\text{now}} < t_{\text{elig}}$
for the entire period before removal.

(2) If detection occurs after $\tau_c$, the agent appears
in discovery for $T_{\text{detect}} - \tau_c$ hours. Since
$\tau_c > 0$, this is strictly less than the exposure window
without cooldown ($T_{\text{detect}}$ hours).
\end{proof}

\subsection{Anti-Sybil Analysis}

\begin{proposition}[Sybil Cost Lower Bound]
\label{prop:sybil}
Under the discovery cooldown, an adversary who seeks to
maintain $k$ agents continuously visible in the discovery
index must register at least:
\begin{equation}
  n_{\text{agents}} \geq k
\end{equation}
agents before the index window opens, at a minimum temporal
cost of $\tau_c = 48\text{ h}$ per batch.  To maintain $k$
agents visible at all times (including during the cooldown
of newly registered replacements), the adversary must
maintain $k$ additional agents in a ``warming up'' pipeline,
requiring at minimum $2k$ total active registrations.
\end{proposition}

\subsection{Direct Resolution Bypass}

The cooldown does not block direct resolution by didHash
(\texttt{GET /v1/agents/:didHash}).  This is a deliberate
design choice: a legitimate agent that has just been
registered should be callable by orchestrators that have
its DID reference (e.g., from a private channel). The
cooldown only prevents \emph{cold discovery}—the case where
a consumer finds the agent by browsing the public index
without prior knowledge of its DID.

% ─────────────────────────────────────────────────────────────
\section{Adversarial Cost Analysis}
\label{sec:cost_analysis}

Table~\ref{tab:cost} summarises the minimum cost for each
attacker profile to defeat each tier.

\begin{table}[t]
\centering
\caption{Minimum attack cost per tier}
\label{tab:cost}
{\small
\begin{tabular}{llp{1.8cm}p{2.0cm}}
\toprule
Tier & Atk. & Attack vector & Cost \\
\midrule
$T_0$ & $\Adv_1$ & Free-text org name
  & $c_{\text{key}} \approx \$0$ \\
$T_1$ & $\Adv_2$ & Typosquat domain
  & $c_{\text{dom}} \approx \$10\text{/yr}$ \\
$T_1$ & $\Adv_3$ & DNS poisoning
  & $p_{\text{dns}}{\cdot}(\text{atk cost})$ \\
$T_2$ & $\Adv_3$ & DNS poison + org-key
  & $p_{\text{dns}}{\cdot}(\text{atk cost}) + \negl(\lambda)$ \\
$T_2$ & Any & Break Ed25519
  & $2^{128}$ ops (infeasible) \\
$T_2$ & Any & Compromise $\sk_{\text{org}}$
  & Key mgmt. dependent \\
\bottomrule
\end{tabular}
}
\end{table}

The progressive cost structure ensures that higher-tier
badges are increasingly expensive to fake. At $T_2$ with
DNSSEC-protected domains, the only practical attack is
compromise of $\sk_{\text{org}}$, which is a key management
problem (HSM usage, key rotation policies) outside the scope
of the registry.

% ─────────────────────────────────────────────────────────────
\section{Comparison to Alternative Approaches}
\label{sec:comparison}

\subsection{X.509 / WebPKI}

Traditional TLS certificate infrastructure assigns
certificates to domain names via Certificate Authorities
(CAs).  DV (Domain Validated) certificates prove domain
control—the same property as our $T_1$ tier, but with
additional CA overhead (>200 KB CA bundle required in the
verifier, CA root store trust, certificate parsing
complexity).  OV/EV certificates include organisation
identity but are rare ($<5\%$ of issued certificates) and
require manual vetting that scales poorly.  We chose our
own well-known binding over X.509 because: (1) it requires
no CA root store in the edge computing environment; (2) it
avoids the CA mis-issuance risk (a rogue CA could issue
a certificate for any domain); (3) the binding is directly
controlled by the domain holder with no third-party
intermediary.

\subsection{DNSSEC}

DNSSEC provides cryptographic authentication of DNS
responses. If universally deployed, it would eliminate the
DNS poisoning risk in our $T_1$ and $T_2$ analyses. In
practice, DNSSEC adoption is approximately 40\% among
top-level domains.  We recommend DNSSEC as a complementary
protection but do not require it, accepting the residual
$p_{\text{dns}}$ risk.

\subsection{Sigstore / OIDC Attestation}

Sigstore \cite{sigstore2021} enables keyless signing via
OIDC tokens from identity providers (GitHub Actions, GCP,
AWS OIDC).  This approach binds a code signing identity to
an OIDC principal (e.g., a GitHub repository), providing
stronger legal-identity guarantees than DNS alone.  We
identify Sigstore integration as the natural evolution of
our $T_2$ tier: instead of a manually managed org keypair,
organisations would sign registration payloads with
ephemeral keys attested by their CI/CD pipeline, providing
verifiable links to GitHub organisations, GCP projects, or
AWS accounts.

% ─────────────────────────────────────────────────────────────
\section{SSRF Mitigations in the Verification Fetch}
\label{sec:ssrf}

The domain verification protocol requires the registry
server to fetch \texttt{https://}$d$\texttt{/.well-known/aroha-did.json}.
This introduces a Server-Side Request Forgery (SSRF) risk
if the attacker controls $d$ and resolves it to an internal
address.

\textbf{Primary defence}: The Aroha registry runs on
Cloudflare Workers, which enforces strict egress policy:
outbound connections to RFC 1918 (private), loopback, and
link-local addresses are blocked at the network layer. A
domain resolving to \texttt{192.168.0.1} or
\texttt{127.0.0.1} produces a connection refusal, not a
successful fetch.

\textbf{Defence in depth}: The registry additionally passes
\texttt{redirect:'manual'} to all well-known fetches, and
rejects any response with \texttt{status >= 300}. This
prevents redirect chains that might reach internal endpoints.
Generic error messages are returned regardless of HTTP status
codes, preventing information leakage about internal topology.

% ─────────────────────────────────────────────────────────────
\section{Implementation}
\label{sec:impl}

\subsection{Database Schema}

The verification fields are stored in the agent registry
database (Cloudflare D1/SQLite):

\begin{verbatim}
ALTER TABLE agent_records ADD COLUMN
  verification_tier TEXT NOT NULL
  DEFAULT 'community';
ALTER TABLE agent_records ADD COLUMN
  domain TEXT;
ALTER TABLE agent_records ADD COLUMN
  domain_verified INTEGER NOT NULL DEFAULT 0;
ALTER TABLE agent_records ADD COLUMN
  discovery_eligible_at INTEGER;
ALTER TABLE agent_records ADD COLUMN
  org_signature TEXT;
\end{verbatim}

The \texttt{discovery\_eligible\_at} column is nullable:
a null value means no cooldown (legacy records); an integer
is the Unix-millisecond timestamp at which the agent
becomes visible. The SQL filter:

\begin{verbatim}
WHERE active = 1
  AND (discovery_eligible_at IS NULL
       OR discovery_eligible_at <= ?)
\end{verbatim}

correctly handles both legacy (null) and cooldown-enabled
records in a single indexed query.

\subsection{Migration Strategy}

Both verification columns are added via \texttt{ALTER TABLE
... ADD COLUMN} with \texttt{DEFAULT} values, which is
zero-downtime in SQLite and PostgreSQL—no locking is
required. Existing records receive \texttt{verification\_tier
= 'community'} and \texttt{discovery\_eligible\_at = NULL}
(no cooldown), preserving backwards compatibility.

% ─────────────────────────────────────────────────────────────
\section{Evaluation}
\label{sec:eval}

\subsection{Verification Latency}

The domain verification endpoint performs one HTTPS fetch
with a 10-second timeout. Across 500 verification requests
against domains with typical TLS round-trip times:

\begin{table}[t]
\centering
\caption{Domain verification latency (ms)}
\label{tab:verify_latency}
\begin{tabular}{lrr}
\toprule
Metric & Domain verify & Org-key register \\
\midrule
Median & $182$ & $211$ \\
P95    & $643$ & $721$ \\
P99    & $1{,}840$ & $2{,}103$ \\
Timeout ($> 10{,}000$) & $0.4\%$ & $0.6\%$ \\
\bottomrule
\end{tabular}
\end{table}

Verification is a one-time operation per agent; the
latency cost is paid once and the tier is permanently
recorded. Discovery queries add zero overhead (tier is
stored in the indexed agent record).

\subsection{Cooldown Effectiveness}

We analysed 6 months of registry data (masked): among $n =
847$ newly registered agents, 14 were flagged by legitimate
organisation owners as impersonation attempts.  Of these:
\begin{itemize}
  \item 9 (64\%) were detected and removed within 48 hours
    of registration, before the cooldown window expired.
    Under the cooldown, these agents had zero discovery
    exposure.
  \item 5 (36\%) were detected after the cooldown expired;
    their discovery exposure was reduced by 48 hours
    compared to the no-cooldown baseline.
\end{itemize}

% ─────────────────────────────────────────────────────────────
\section{Limitations and Future Work}
\label{sec:limits}

The current $T_2$ tier does not achieve legal-identity
verification—only domain-key binding with DNS trust.
Sigstore integration (\S\ref{sec:comparison}) is the primary
future direction. Additionally, the well-known binding
requires the publisher to maintain a publicly accessible
HTTPS server; this excludes agents published by entities
that do not control public-facing web infrastructure (e.g.,
internal enterprise agents). Enterprise VPN-accessible
well-known endpoints or challenge-response mechanisms are
candidate solutions.

% ─────────────────────────────────────────────────────────────
\section{Conclusion}
\label{sec:conclusion}

We have presented a three-tier publisher verification
architecture for decentralised AI agent registries. The
architecture progressively raises the cost and cryptographic
strength of publisher identity claims: from free-text
metadata ($T_0$) to DNS-controlled well-known binding ($T_1$)
to Ed25519 org-keypair signing ($T_2$). We formalised the
threat model with three attacker profiles, proved security
properties for each tier, and showed that the discovery
cooldown reduces impersonation exposure for the majority
of detection scenarios. Comparison with X.509, DNSSEC, and
Sigstore identifies a clear evolutionary path toward
stronger legal-identity guarantees. The system is deployed
in production on the Aroha agent registry.

\bibliographystyle{IEEEtran}
\bibliography{aroha}

% ─────────────────────────────────────────────────────────────
\appendix

\section{Well-Known File Format}
\label{app:wellknown}

\textbf{Domain tier} (\texttt{/.well-known/aroha-did.json}):
\begin{verbatim}
{
  "did": "did:aroha:studio:0x<keccak256-of-did>"
}
\end{verbatim}

\textbf{Org-key tier} (\texttt{/.well-known/aroha-org.json}):
\begin{verbatim}
{
  "orgPublicKeyB64": "<base64-encoded-32-byte-ed25519-pubkey>"
}
\end{verbatim}

Both files must be served over HTTPS with a valid TLS
certificate. HTTP-only domains fail verification.

\section{Registration Signature Example}
\label{app:sig_example}

The canonical signed payload for org-signed registration is:

\begin{verbatim}
{
  "action": "register-agent",
  "did": "did:aroha:studio:0xabc...",
  "manifestCID": "bafyreib...",
  "publicKeyB64": "AAEC...",
  "registeredAt": 1735689600000
}
\end{verbatim}

Note: keys are sorted ASCII-lexicographically per
RFC 8785 \cite{jsoncanon}. Any deviation in key order or
whitespace produces a different byte sequence and an invalid
signature. The registry uses the same canonical encoder for
both the signing check and the message construction, ensuring
bit-for-bit compatibility.

\end{document}
